During the MCMC exploration, we occasionally encounter runs that terminate abnormally or converge unrealistically fast. These cases are characterized by near-zero acceptance fractions, undefined or NaN autocorrelation times, and extremely short execution times. Such behavior indicates that the ensemble of walkers fails to effectively explore the parameter space and instead becomes trapped at invalid or zero-probability regions, often due to overly restrictive priors, numerical instabilities in the likelihood evaluation, or strong parameter degeneracies that collapse the walkers onto a lower-dimensional subspace. In contrast, successful MCMC runs typically exhibit longer execution times, acceptance fractions within the expected range, and well-defined autocorrelation times, reflecting efficient sampling of the posterior distribution. We therefore interpret runs with vanishing acceptance rates and undefined autocorrelation estimates as non-converged and exclude them from further analysis, retaining only those chains that demonstrate stable sampling behavior and statistically meaningful convergence diagnostics.